\documentclass[10pt]{article}

\usepackage[margin=1in]{geometry}
\usepackage{booktabs}
\usepackage{hyperref}
\usepackage{url}
\usepackage{xspace}

\newcommand{\tool}{Bayesilisk\xspace}

\title{Beyond E2E Scripts: Using LLM-Proposed Scenarios Without Letting the LLM Be the Oracle}
\author{Draft manuscript}
\date{May 2026}

\begin{document}
\maketitle

\begin{abstract}
Large language models can help testers explore application states that are not
covered by ordinary end-to-end scripts, but they are unsafe as test oracles.
They can suggest plausible scenarios, misread application semantics, and
overstate weak evidence. \tool separates these roles. App-specific connectors
publish observable context, available actions, and deterministic rules.
\tool expands the supplied scenario space, ranks probes, and verifies only
observed evidence returned by the connector. This draft reports an initial
Cal.com case study at commit
\texttt{180ede28f0bddf2738933a6e60a8e80f6116d7da}: connector-supplied
route facts and an action graph produced seven Bayesilisk probe proposals. Six
local Playwright executions reproduced a mismatch where unknown or stale
\texttt{rescheduleUid} values opened booking flows with semantic status
\texttt{200} although the supplied invariant expected \texttt{404}; a seventh
bounded workflow proposal created a booking, cancelled it, and replayed the
cancelled booking UID through the public booking route, again observing
\texttt{200} where the supplied invariant expected \texttt{409}. One reported
finding already has an upstream human-authored fix pull request. The current
artifact set keeps only this clean seven-finding rerun; baseline comparisons
should be recreated as fresh runs when needed.
\end{abstract}

\section{Problem}

E2E suites are precise but narrow. LLM-driven agents are broad but unreliable:
they can drive a browser, invent scenarios, and then incorrectly decide whether
the behavior is acceptable. The core research claim behind \tool is that these
jobs should be separated. An LLM or coding agent may help propose scenario
spaces, but deterministic verifier rules must decide pass or fail.

The intended workflow is:

\begin{enumerate}
  \item A connector extracts app-specific context: routes, actors, identifiers,
        available actions, nearby tests, and expected behavior.
  \item \tool expands only the supplied rule space and produces probe proposals.
  \item The connector executes those proposals through the real app or its test
        harness.
  \item \tool verifies returned evidence against deterministic invariants and
        emits issue-ready findings.
\end{enumerate}

The connector is app-specific. The Bayesilisk core is not. This boundary is the
main architectural point: app teams write connectors; \tool provides the
generic expansion, ranking, and verification loop.

\section{Cal.com Case Study}

We evaluated the current connector architecture on a local Cal.com checkout.
The tested repository and revision were:

\begin{itemize}
  \item Repository: \url{https://github.com/calcom/cal.com}
  \item Commit: \texttt{180ede28f0bddf2738933a6e60a8e80f6116d7da}
  \item Commit date: \texttt{2026-05-14 19:30:21 +0000}
  \item Commit subject: font-stack fallback fix for non-Latin script support
        (\texttt{\#29346})
  \item Bayesilisk seed: \texttt{150}
\end{itemize}

The checked-in Bayesilisk artifacts are in \texttt{examples/calcom/}. They
include the connector, source contexts, generated proposals, consolidated
execution context, app-only JSON report, app-only Markdown report, issue
payloads, and the two Playwright run logs for the clean rerun.

\begin{table}[h]
\centering
\begin{tabular}{ll}
\toprule
Artifact & Path \\
\midrule
Connector code & \texttt{examples/calcom/bayesilisk-probes.e2e.ts} \\
Source context & \texttt{examples/calcom/source-context.json} \\
Generated proposals & \texttt{examples/calcom/generated-proposals.json} \\
Sequence context & \texttt{examples/calcom/sequence-source-context.json} \\
Generated sequence & \texttt{examples/calcom/generated-sequence-proposals.json} \\
Execution evidence & \texttt{examples/calcom/execution-context.json} \\
Bayesilisk report & \texttt{examples/calcom/reports/report.json} \\
Issue payloads & \texttt{examples/calcom/reports/issue-payloads.json} \\
Upstream outcomes & \texttt{examples/calcom/upstream-outcomes.md} \\
\bottomrule
\end{tabular}
\caption{Checked-in Cal.com artifacts.}
\end{table}

The connector supplied source facts for Cal.com routes and an action graph for
one workflow. \tool generated six route probes and one workflow-sequence probe.
The route probes mutate unknown and stale \texttt{rescheduleUid} values for
public, private, and dynamic booking routes. The workflow probe creates a
booking, cancels it, and replays the cancelled booking UID through the public
booking route. The connector executed all seven proposals with Playwright
against local fixtures.

\begin{table}[h]
\centering
\begin{tabular}{llll}
\toprule
Area & Mutation & Expected & Observed \\
\midrule
Public booking page & unknown id & 404 & 200 \\
Public booking page & stale id & 404 & 200 \\
Private booking link & unknown id & 404 & 200 \\
Private booking link & stale id & 404 & 200 \\
Dynamic booking page & unknown id & 404 & 200 \\
Dynamic booking page & stale id & 404 & 200 \\
Cancelled booking replay & cancelled uid & 409 & 200 \\
\bottomrule
\end{tabular}
\caption{Verified Cal.com connector observations.}
\end{table}

The result is useful because the pass/fail verdict did not come from a model.
The connector returned concrete observed facts. \tool compared
\texttt{expectedStatus} and \texttt{observedStatus}, classified the seven
observations as \texttt{breakage.context-observed}, and produced issue payloads.
Two upstream issues were opened from this evidence: one for public
\texttt{rescheduleUid} replay~\cite{calcomIssue} and one for cancelled booking
UID replay~\cite{calcomCancelledIssue}. The first issue already has a
human-authored fix pull request~\cite{calcomFixPr}. This is stronger validation
than issue closure alone: an issue can be closed as erroneous, but a targeted
fix PR indicates that a human contributor found the report actionable.

\section{Upstream Response}

The upstream response is important because the claim is not only that
\tool can emit plausible reports. The reports need to be useful to humans
working on the target system.

\begin{table}[h]
\centering
\begin{tabular}{llll}
\toprule
Finding group & Issue & Human response & Status at capture \\
\midrule
Public unknown/stale \texttt{rescheduleUid} & \#29399 & Fix PR \#29400 & Open \\
Private-link \texttt{rescheduleUid} & \#29399 comment & Folded into thread & Open \\
Cancelled booking UID replay & \#29407 & Newly reported & Open \\
\bottomrule
\end{tabular}
\caption{Upstream response to Bayesilisk-generated Cal.com findings.}
\end{table}

The private-link finding is tracked as a related comment rather than a separate
issue to avoid duplicating the existing \texttt{rescheduleUid} thread. The
cancelled-booking workflow finding is separate because it comes from a different
state transition: a real booking is created, cancelled, and then replayed as
public reschedule context.

\section{Evidence Still Needed}

The Cal.com result demonstrates the connector split and a real route/state
mutation finding. It does not yet demonstrate the full paper claim. The
remaining evidence needs are fresh baseline comparisons and additional app
artifacts beyond Cal.com.

\begin{table}[h]
\centering
\begin{tabular}{llll}
\toprule
Workflow & Observed scope & Failures & Notes \\
\midrule
Bayesilisk proposal run & 6 proposals & 6 & Unknown and stale variants \\
Bayesilisk sequence run & 1 sequence & 1 & Cancelled UID replay \\
\bottomrule
\end{tabular}
\caption{Current clean Cal.com rerun.}
\end{table}

Fresh baselines matter because the argument is not merely that browser
automation can find a failing route. The intended claim is stronger: generic
context-to-probe expansion plus deterministic verification can recover useful
failures while avoiding the oracle mistakes of plain LLM agents.

\section{Discussion}

The current evidence supports a narrow claim: given connector-supplied source
facts and rules, \tool can expand a small scenario space, execute it through an
app-specific connector, and verify evidence without letting the LLM decide the
outcome. This is already different from a plain Playwright script because the
connector does not enumerate every final generated probe by hand; it exposes
the local state and rule space, and \tool produces the concrete mutations.

The remaining work is empirical. The paper needs more app artifacts and
architecture work before it can claim generality for longer workflows. The new
sequence proposer is intentionally bounded: connectors declare actions and
state facts, while \tool generates short action sequences and verifies returned
evidence. The immediate target should be a concise artifact-backed paper rather
than a broad methodology claim.

\section{Conclusion}

\tool treats LLMs as scenario proposers, not as test oracles. The Cal.com case
study shows the architecture working on a real app checkout with reproducible
connector evidence and issue-ready findings. The next version of this
manuscript should add a full second app artifact and recreate baseline
comparisons against existing E2E tests, Playwright-only probing, weak model
proposals, and browser-driving LLM agents.

\bibliographystyle{plain}
\bibliography{references}

\end{document}
